% !TeX root = ../main.tex

\chapter{結論與未來工作}
\label{chap:conclusion}

\section{研究總結}
\label{sec:summary}

本論文針對自免疫螢光染色之皮膚切片影像自動量測表皮內神經纖維密度 (IENFD) 之實務問題，指出於 IENF 片段切分至符合 EFNS 計數規則之完整纖維拓樸之間，存在一個\textbf{中介重建步驟之空缺}：片段切分結果本即由多段不連續之纖維微粒組成，單純骨架化並無法得到可直接計數之拓樸。為填補此空缺，本研究將其形式化為\textbf{像素成本圖}上之多源最短路徑與圖論重建問題，並提出\textbf{標註擴張連結} (AEL)：以影像之纖維響應作為成本圖，將各標註微粒視為獨立來源向外擴張，於相鄰微粒之最小成本相遇處建立連結，再經拓樸過濾與骨架化，產生可供有效神經纖維計數之像素級拓樸圖。

相較於既有基準方法，AEL 具備以下三項主要優勢：

\begin{itemize}
  \item \textbf{影像感知之連結準則}：連結路徑沿影像上實際存在之纖維線索延伸，而非如純距離連結般以直線捷徑跨越訊號空隙，從而抑制偏離纖維之多餘連線。
  \item \textbf{兩階段拓樸過濾}：閾值裁剪抑制偽連線、最小生成森林抑制冗餘連線，兩者皆於元件層上操作，計算成本低且可解釋。
  \item \textbf{可稽核之完整流程}：輸出為可視化之拓樸圖而非純量結果，每一次被計為神經纖維之線段皆對應影像上具體之像素路徑，便於臨床複核。
\end{itemize}

於本研究所使用之資料集上，AEL 於有效神經纖維計數之三項指標（$N_{\text{MAE}}$、$N_{\text{RMSE}}$、$\rho$）與拓樸品質之三項指標（$\text{Dice}_{\text{cl}}$、HD95、$\text{HD}_{\text{avg}}$）上，均優於 MSF 連結之基準方法；結構性消融亦顯示前處理、成本設計與連結策略各模組對最終拓樸品質皆有正向且多數顯著之貢獻。參數分析方面，前處理與線狀結構濾波階段之五個超參數，其預設值係經與鄰近取值逐一比較後選定而非任意指定，連結階段之邊裁剪閾值 $\tau$ 對重建品質之影響最為顯著。

除演算法本身外，本研究亦實作一套\textbf{視覺化神經纖維分析工具}（第~\ref{chap:tool}~章），將 AEL 之重建結果以可互動之圖層疊合於原始影像之上，使專家得以逐條檢視重建之纖維，並於介面上直接增刪節點與邊以修正重建、即時重新計數。此工具將 AEL 之可稽核性落實為「自動重建—人工複核—修正」之半自動工作流程，為本研究除演算法外之另一項貢獻。

IENFD 之自動量測，長期受制於「片段化之纖維偵測」與「臨床所需之完整纖維計數」之間的落差。本論文之核心貢獻，在於指認並填補了兩者之間的\textbf{中介重建空缺}，說明此一步驟並非可有可無之後處理，而是連接語意分割與臨床計數規則之關鍵環節。期望本研究對此空缺之形式化與初步解法，能成為後續「分割—重建—計數」整合式研究之起點，最終推動小纖維神經病變之客觀量測在臨床上之廣泛採用。

\section{研究限制}
\label{sec:limitations}

本研究之結果受限於以下若干因素：

\begin{enumerate}
  \item \textbf{資料集規模與多樣性}：本研究所使用之資料集來源單一，切片厚度與染色協定一致，遷移性仍待驗證。
  \item \textbf{對上游分割品質之依賴}：AEL 僅能以「既有標註」為擴張來源，若某纖維於整段長度上皆未被捕捉，AEL 將無法自無中生有地加入該纖維。
  \item \textbf{表皮遮罩之精確性}：有效神經纖維之判定直接依賴 $M$，在手工標註誤差或自動表皮分割失準時會造成計數偏移。
  \item \textbf{二維切片之限制}：本研究所處理之影像為單一垂直切片之二維投影；真實纖維於三維空間中可能繞出又進入同一切片，於二維上被計為多次有效神經纖維。此問題需透過真正的三維影像處理方能徹底解決。
\end{enumerate}

\section{未來工作}
\label{sec:future}

基於上述結果與限制，本研究提出以下幾項具體之延伸方向：

\subsection{與語意分割之聯合訓練}

現行 AEL 將語意分割之輸出視為固定輸入。若將 AEL 之重建結果反饋至分割模型之訓練目標，使分割模型偏好輸出「易於被 AEL 補足之纖維微粒」而非「形狀完整但本質錯誤之大塊」，有機會形成一個\textbf{分割—重建雙向優化}之迴圈。一種具體之實作是：將 AEL 輸出之重建骨架與專家有效神經纖維地面真值之差異回傳為分割訓練之損失函數。

\subsection{三維切片之延伸}

隨著組織透明化 (tissue clearing) 與共聚焦／螢光三維掃描技術之成熟，研究級資料已可取得三維之神經纖維影像。AEL 之架構自然可延伸至三維：僅需將 8-連通鄰域改為 26-連通，並於 3D Hessian 矩陣上計算 Sato 線狀結構響應即可。在三維中，有效神經纖維 DEJ 之判定亦需改為與一個 DEJ 曲面之相交關係，而非與 2D 遮罩之相交。

\subsection{可學習之成本函數}

本研究之成本函數為手動設計之固定映射。若以小型卷積網路將影像映射為成本圖、並以 AEL 之下游計數誤差為監督訊號進行反向傳播，則有機會進一步提升成本圖對 DEJ 附近染色中斷之適應性。

\subsection{不確定性量化}

對於臨床決策，除了 IENFD 之點估計外，其不確定性亦具重要意義。未來可為每一條被計為有效神經纖維之線段輸出對應之信心分數，並進一步量化整張切片 IENFD 之信賴區間，以輔助臨床判讀。
